\chapter{Gemelo termodinámico del CFM56-7B26}
\label{cap:modelo}

Este capítulo documenta el trabajo completado de la fase F1: el modelo de prestaciones del
CFM56-7B26 en pyCycle~\cite{Hendricks2019}, su punto de diseño (WP1.1) y su calibración fuera de
diseño contra datos medidos certificados (WP1.2). Todas las cifras de resultados provienen del
script de generación de evidencias (\cref{sec:reproducibilidad}).

\section{Estrategia: adaptar, no construir desde cero}
\label{sec:morph}

pyCycle distribuye un ejemplo validado de turbofan de alta derivación de doble rotor y flujo
separado con refrigeración de turbinas (\codigo{high\_bypass\_turbofan.py}), con el grafo de
elementos, la red de sangrados, los balances de ejes y el solver fuera de diseño ya operativos.
Construir el CFM56 desde un fichero en blanco obliga a re-derivar todo ese cableado y suele
morir en divergencias del método de Newton. En su lugar se adopta un protocolo de
\emph{morphing}: el ejemplo se toma como estado conocido-bueno (y se conserva en el repositorio
como referencia de regresión) y se re-orienta hacia el CFM56-7B26 cambiando \textbf{como máximo
tres parámetros por paso}, exigiendo convergencia en cada paso antes del siguiente. Cada paso
convergido es un \emph{commit}; si un paso diverge, se bisecta el cambio.

Las decisiones de diseño se fijaron por escrito \emph{antes} de codificar
(\codigo{docs/f1-model-spec.md}): variante -7B26; termodinámica tabular con verificación cruzada
CEA; punto de diseño en crucero M\,0,78 / \SI{35000}{ft} / ISA; mapas genéricos de pyCycle
escalados (no existen mapas públicos del CFM56) con extrapolación controlada; potencia fijada
por N1 corregida en fuera de diseño (el -7B es un motor regulado por N1); unidades nativas de
pyCycle en el interior del modelo y SI solo en la frontera; y una tabla de modos de fallo del
solver con su contramedida prevista, que resultó profética (\cref{sec:idle}).

\section{Punto de diseño (WP1.1)}
\label{sec:diseno}

El grafo del ciclo comprende: condiciones de vuelo, entrada, fan, divisor de flujo,
\emph{booster}, HPC con tres puertos de sangrado (refrigeración de vano y álabe de HPT y
sangrado de cliente), sangrado post-compresor para la LPT, cámara, HPT y LPT refrigeradas,
conductos con pérdidas de presión, toberas convergentes de núcleo y derivación, dos ejes y
extracción de potencia de accesorios (\SI{250}{hp} en el eje de alta). En el punto de diseño,
cuatro balances implícitos determinan el ciclo: el gasto $W$ ajusta el empuje objetivo; el
dosado FAR ajusta $T_{t4}$; y las relaciones de expansión de HPT y LPT anulan la potencia neta
de cada eje.

\begin{table}[htbp]
  \centering\small
  \caption{Resultados del punto de diseño (M\,0,78, \SI{35000}{ft}, ISA). Tabla autogenerada
  desde el modelo.}
  \label{tab:diseno}
  % AUTOGENERADO por scripts/make_memoria_assets.py — NO editar a mano
\begin{tabular}{ll}
\toprule
Magnitud & Valor en el punto de diseño \\
\midrule
  Empuje neto $F_n$ & \SI{24.38}{\kilo\newton} (\num{5480} lbf) \\
  TSFC & \num{0.6214} lb/(lbf$\cdot$h) \\
  Relación de derivación (BPR) & \num{5.10} \\
  Relación de compresión global (OPR) & \num{30.03} \\
  $T_{t4}$ & \SI{1587}{\kelvin} (\num{2857} °R) \\
  Gasto másico $W$ & \SI{142.0}{\kilo\gram\per\second} \\
  Dosado FAR & \num{0.0251} \\
  Relación de expansión HPT & \num{3.586} \\
  Relación de expansión LPT & \num{4.288} \\
  Fracción total de refrigeración/sangrado & \num{0.239} \\
\bottomrule
\end{tabular}

\end{table}

\begin{table}[htbp]
  \centering\small
  \caption{Condiciones totales por estación en el punto de diseño. Tabla autogenerada.}
  \label{tab:estaciones}
  % AUTOGENERADO por scripts/make_memoria_assets.py — NO editar a mano
\begin{tabular}{llrr}
\toprule
Estación & Posición & $P_t$ [bar] & $T_t$ [K] \\
\midrule
  2 & Entrada fan & 0.358 & 245.5 \\
  21 & Salida fan & 0.603 & 289.5 \\
  25 & Entrada HPC & 1.149 & 354.2 \\
  3 & Salida HPC & 10.739 & 703.9 \\
  4 & Salida cámara & 10.159 & 1587.2 \\
  45 & Salida HPT & 2.833 & 1141.6 \\
  5 & Salida LPT & 0.657 & 806.3 \\
\bottomrule
\end{tabular}

\end{table}

\begin{figure}[htbp]
  \centering
  \includegraphics[width=0.72\textwidth]{estaciones_diseno}
  \caption{Evolución de presión y temperatura totales a lo largo del camino de gas en el punto
  de diseño. La compresión culmina en la estación 3 ($\sim$\SI{11}{bar},
  \cref{tab:estaciones}), la adición de calor en la 4 y la extracción de trabajo en 45 y 5.}
  \label{fig:estaciones}
\end{figure}

El punto de diseño converge con residuo de Newton inferior a $10^{-6}$ y satisface las cinco
ventanas de aceptación definidas a priori: BPR $5{,}10\pm0{,}2$ (objetivo EEDB), OPR en ventana,
$T_{t4}$ en banda plausible, fracción total de refrigeración y sangrado en 18--25\,\% del flujo
de núcleo, y empuje exacto al objetivo (\cref{tab:diseno}). El TSFC resultante
($\approx0{,}62$~lb/(lbf$\cdot$h)) queda a menos del 2\,\% del valor de crucero citado
públicamente para este motor, sin haberse usado como objetivo de ajuste.

\section{Calibración contra datos medidos (WP1.2)}
\label{sec:calibracion}

\subsection{Fuentes de datos y filosofía}

La calibración fuera de diseño usa exclusivamente \textbf{datos medidos y certificados}
(\cref{tab:tcds}): la TCDS aporta límites y empujes nominales, y el banco de emisiones OACI
aporta los flujos de combustible medidos en banco a nivel del mar en los cuatro reglajes del
ciclo aterrizaje-despegue (100, 85, 30 y 7\,\% del empuje nominal $F_{00}$), además de la
relación de presiones del motor al régimen nominal (27{,}61). Estos cuatro puntos ---ignorados
por la mayoría de trabajos académicos--- constituyen anclas de calibración gratuitas, trazables
y de calidad de certificación.

Los puntos ancla se resuelven \emph{a empuje impuesto} (el dosado equilibra
$F_n = \mathrm{PC}\cdot F_{00}$): así el flujo de combustible del modelo es una
\textbf{predicción} genuina que se confronta con la medida del EEDB, no un ajuste.

\subsection{Iteraciones de calibración}

La primera ejecución (modelo recién adaptado del ejemplo, sin calibrar a los datos SLS) predijo
el flujo de combustible de despegue con un error de $-2{,}1$\,\%, pero presentó dos defectos
cuantitativos:

\begin{enumerate}
  \item \textbf{N2 al empuje nominal por encima de la línea roja}: \SI{15311}{rpm} frente al
    límite certificado de \SI{15183}{rpm}. Causa: la referencia de velocidad con la que se
    escala el mapa genérico del HPC. Corrección: referencia de diseño del eje de alta de
    \SI{14324}{} a \SI{13940}{rpm}, con lo que el régimen nominal queda en \SI{14915}{rpm}
    ($\approx103$\,\%), dentro del límite.
  \item \textbf{OPR de despegue $28{,}99$ frente al valor medido $27{,}61$} ($+5{,}0$\,\%).
    Corrección: relación de compresión del HPC en el punto de diseño de $9{,}81$ a $9{,}35$; el
    dato medido prevalece sobre la estimación derivada inicial, y el OPR de crucero pasa a ser
    una consecuencia ($\approx30$).
\end{enumerate}

\begin{table}[htbp]
  \centering\small
  \caption{Anclas SLS frente a las medidas del banco OACI tras la calibración. $W_f$~mod.\ es la
  predicción del modelo a empuje impuesto. Tabla autogenerada.}
  \label{tab:anclas}
  % AUTOGENERADO por scripts/make_memoria_assets.py — NO editar a mano
\begin{tabular}{lrrrrrrc}
\toprule
Ancla & \%$F_{00}$ & $F_n$ [kN] & $W_f$ mod. & $W_f$ EEDB & error [\%] & tol. [\%] & gate \\
\midrule
  A2 Takeoff & 100 & 117.0 & 1.204 & 1.221 & -1.40 & \num{3.0} & \checkmark \\
  A3 Climbout & 85 & 99.4 & 0.977 & 0.999 & -2.23 & \num{3.0} & \checkmark \\
  A4 Approach & 30 & 35.1 & 0.328 & 0.338 & -3.08 & \num{5.0} & \checkmark \\
  A5 Idle & 7 & 8.2 & 0.125 & 0.113 & +10.57 & \num{5.0} & --- \\
\bottomrule
\end{tabular}

\end{table}

\begin{figure}[htbp]
  \centering
  \includegraphics[width=0.7\textwidth]{errores_anclas}
  \caption{Error de la predicción de flujo de combustible en cada ancla frente a la medida del
  EEDB. Las bandas grises son las tolerancias del contrato ($\pm3$\,\% en despegue y
  \emph{climb-out}, $\pm5$\,\% en aproximación e idle); el punto hueco (idle) converge pero
  queda fuera de tolerancia y se reporta sin bloquear el hito (\cref{sec:idle}).}
  \label{fig:errores-anclas}
\end{figure}

\begin{figure}[htbp]
  \centering
  \includegraphics[width=0.7\textwidth]{linea_trabajo_sls}
  \caption{Línea de trabajo del modelo a nivel del mar estático (barrido continuo de reglaje de
  empuje) frente a los cuatro puntos medidos del banco OACI. La curvatura a baja potencia
  refleja el deterioro del rendimiento con mapas genéricos lejos del punto de escalado.}
  \label{fig:linea-trabajo}
\end{figure}

Tras ambas correcciones (\cref{tab:anclas,fig:errores-anclas,fig:linea-trabajo}), las tres
anclas con \emph{gate} cumplen sus tolerancias: despegue $-1{,}40$\,\% de error en $W_f$ con OPR
a $+0{,}3$\,\% del medido y ambos regímenes dentro de las líneas rojas; \emph{climb-out}
$-2{,}23$\,\%; aproximación $-3{,}08$\,\%.

\subsection{Convergencia en idle: diagnóstico y solución}
\label{sec:idle}

El punto de idle (7\,\% de $F_{00}$, estático) diverge sistemáticamente desde arranque en frío:
con relaciones de expansión de tobera próximas a la unidad, el subsolver de presión estática de
las toberas rebota contra su cota inferior y el Newton global no encuentra dirección de
descenso. Peor aún, el intento fallido deja los estados internos del punto (estaciones de flujo,
posiciones de mapa) corrompidos, de modo que re-sembrar solo los balances no lo rescata: fue
necesario que el punto de idle \emph{nunca} arranque en frío en su reglaje final. La solución
---prevista como contramedida genérica en la especificación--- es una \textbf{continuación en
potencia}: el punto arranca al 30\,\% de $F_{00}$ (donde converge de frío, igual que el ancla de
aproximación) y desciende en caliente por la secuencia $0{,}30 \rightarrow 0{,}20 \rightarrow
0{,}14 \rightarrow 0{,}10 \rightarrow 0{,}08 \rightarrow 0{,}07$, re-resolviendo el problema en
cada paso. Esta lógica queda encapsulada en la librería (\codigo{solve\_anchors()}) y cubierta
por las pruebas.

Con la continuación, el idle converge y predice $W_f$ con $+10{,}6$\,\% de error. Es territorio
de extrapolación profunda de los mapas genéricos (la forma velocidad-flujo a muy baja potencia
no es la del compresor real), una limitación conocida del enfoque; el ancla se \textbf{reporta
pero no bloquea} el hito H1, decisión documentada en el contrato de calibración
(\codigo{gated:~false}).

\section{Limitaciones declaradas}
\label{sec:limitaciones}

El modelo es \emph{representativo}, no certificable, y sus límites se declaran para su cita en
los capítulos de resultados: (i) sin transitorios ---solo estados estacionarios---; (ii)
geometría variable (VBV/VSV) implícita en los mapas, no modelada; (iii) sin correcciones de
Reynolds ni humedad; (iv) las etiquetas absolutas de N1/N2 son aproximadas por construcción
(formas de mapa genéricas), aunque respetan las líneas rojas certificadas en el punto nominal;
(v) la EGT del modelo es la temperatura de salida de la HPT (estación 45), un \emph{proxy}
coherente ---idéntico para la verdad-terreno y para ambos pipelines de EHM--- de la T49.5
certificada, cuyo valor mostrado además incorpora las funciones \emph{shunt} y \emph{trim}
descritas en \cref{sec:motor}; y (vi) el error de idle citado arriba. Ninguna de estas
limitaciones afecta a la validez de la comparación IA-contra-tradicional, que se apoya en la
estructura de sensibilidades del modelo (la ICM), no en sus valores absolutos.

\section{Verificación automatizada}

El estado del capítulo queda blindado por pruebas: convergencia y ventanas del punto de diseño;
convergencia de las cuatro anclas con la continuación de idle; tolerancias de $W_f$ y OPR en las
anclas con \emph{gate}; líneas rojas de N1/N2 al empuje nominal; y la prueba de humo del entorno
pyCycle. Nueve pruebas en verde en integración continua en el momento de escribir este
capítulo.
